\section{Volitional Energy}

\begin{definition}[Volitional energy]
Volitional energy $\Will(t)$ is the limited capacity available at time $t$ to
interrupt the current behavioural frame and initiate stabilizing action toward
a target frame.
\end{definition}

\begin{definition}[Leverage]
The effective leverage of free will at time $t$ is
\[
  \Leverage(t) = \frac{\Will(t)}{\DV(t)}.
\]
\end{definition}

\begin{theorem}[Decision-window principle]
Free will has maximum strategic effect not when effort is most intense, but
when the activation barrier is locally minimal relative to available
volitional energy.
\end{theorem}

\begin{proof}[Proof sketch]
If leverage is represented by $\Will(t)/\DV(t)$, then leverage rises when
$\Will(t)$ increases or $\DV(t)$ decreases. In ordinary life, large increases
in available will are unreliable, but temporary drops in barrier are common:
after waking, after finishing a task, upon entering a new environment, or
immediately after a sharp emotional update. Therefore the most reliable
strategy is to act during barrier troughs rather than waiting for heroic
motivation.
\end{proof}
